% Plano de Trabalho de TCC01 - Engenharia de Computação UNIFEI
% Formato A4, margens superior 1,5 cm; inferior 2,5 cm; esquerda e direita 2,0 cm
% Parágrafos: 0 pt (Antes), 6 pt (Depois) e 1,5 linha (Entre linhas)
% Fonte do corpo: Arial 10pt; Títulos: Arial 12pt negrito
% Limite estrito UNIFEI: Até 04 (quatro) páginas excetuando-se a capa e incluindo referências

\documentclass[a4paper,10pt]{article}

\usepackage{cmap}
\usepackage[english,brazil]{babel}
\usepackage{framed}
\usepackage{hyperref}
\usepackage{amsmath}
\usepackage{graphicx}
\usepackage[colorinlistoftodos]{todonotes}
\usepackage{wrapfig}
\usepackage{listings}
\usepackage{color}
\usepackage{indentfirst}
\usepackage{times}
\usepackage{textcomp}
\usepackage{pgfgantt}
\usepackage{booktabs}
\usepackage{colortbl}
\usepackage{ifluatex,ifxetex}
\usepackage{enumitem}
\setlist{nosep, leftmargin=*, topsep=2pt, partopsep=0pt}

% Suporte híbrido transparente para Arial (XeLaTeX no Windows / helvet no pdfLaTeX no Linux)
\ifxetex
  \usepackage{fontspec}
  \setmainfont{Arial}
\else\ifluatex
  \usepackage[utf8]{inputenc}
  \usepackage[T1]{fontenc}
  \usepackage[scaled]{helvet}
  \renewcommand{\familydefault}{\sfdefault}
\else
  \usepackage[utf8]{inputenc}
  \usepackage[T1]{fontenc}
  \usepackage[scaled]{helvet}
  \renewcommand{\familydefault}{\sfdefault}
\fi\fi

\usepackage[top=15mm,bottom=18mm,left=20mm,right=20mm]{geometry}
\usepackage{setspace}\setstretch{1.16} % espaçamento ajustado para conformidade estrita de 4 páginas
\usepackage{titlesec}
\titleformat*{\section}{\large\bfseries}
\titleformat*{\subsection}{\normalsize\bfseries}
\titleformat*{\subsubsection}{\normalsize\bfseries}
\titlespacing*{\section}{0pt}{5pt}{2pt}
\titlespacing*{\subsection}{0pt}{3pt}{1pt}
\setlength{\parskip}{2.2pt}
\setlength{\parindent}{0.8cm}

\newif\ifblackandwhite
\blackandwhitetrue

\usepackage{etoolbox}
\usepackage{longtable}
\usepackage{pdflscape}

\newcommand{\myrowcolour}{\rowcolor[gray]{0.925}}

\ifblackandwhite
  \newcommand{\cheading}[2]{\textbf{#1\hfill #2}}
  \newcommand{\highest}[1]{\textbf{#1}}
\else
  \newcommand{\cheading}[2]{\textcolor{Maroon}{\textbf{#1\hfill #2}}}
  \newcommand{\highest}[1]{\textcolor{Maroon}{\textbf{#1}}}
\fi

\definecolor{mygray}{rgb}{0.4,0.4,0.4}
\definecolor{mygreen}{rgb}{0,0.8,0.6}
\definecolor{myorange}{rgb}{1.0,0.4,0}

\newcommand{\HRule}{\rule{\linewidth}{0.5mm}}

\begin{document}

\begin{titlepage}
\begin{center}

% Logo Institucional UNIFEI
\includegraphics[width=0.28\textwidth]{img/unifei_logo_color.png}~\\[0.8cm]

\textsc{\Large Universidade Federal de Itajubá}\\[0.3cm]
\textsc{\large Instituto de Engenharia de Sistemas e Tecnologias da Informação -- IESTI}\\[0.3cm]
\textsc{\large Curso de Engenharia de Computação}\\[1.8cm]

{\LARGE \bfseries Integração de Modelos de Linguagem e Automação de Fluxos para Apoio à Tomada de Decisão Baseada em Dados}\\[1.8cm]

% Identificação do Edital / Finalidade
\HRule \\[0.3cm]
{\large \bfseries Plano de Trabalho de TCC01 apresentado para análise e deliberação da Coordenação de Trabalho de Conclusão de Curso de Engenharia de Computação}
\HRule \\[1.5cm]

% Identificação do Discente
\large\textbf{Discente:}\\
Gabriel Gonçalves Sampaio -- Engenharia de Computação\\[0.8cm]

% Identificação da Orientadora
\large\textbf{Orientação:}\\
Profa. Dra. Barbara Pimenta Caetano -- Orientadora\\[1.5cm]

\vfill

% Data e Localidade
{\large Itajubá -- MG \\ \today}

\end{center}
\end{titlepage}

\newpage

\section{Introdução}
\label{sc:introducao}

Nos últimos anos, mesmo com a queda nos índices de desmatamento no Brasil, a ocorrência de queimadas não diminuiu — impulsionados por secas severas, os incêndios atingiram níveis recordes na Amazônia, Cerrado e Pantanal. O avanço do fogo afeta diversas frentes: no âmbito socioambiental, mais de 160 Terras Indígenas registraram focos de calor, incluindo o incêndio que destruiu 38 moradias na Aldeia Canoanã, na Ilha do Bananal (TO) \cite{g1_canoana_2026}; na saúde e logística, a fumaça densa elevou atendimentos respiratórios no Sistema Único de Saúde (SUS) e interditou rodovias federais; e no setor elétrico, o calor e a fuligem sob linhas de transmissão do Sistema Interligado Nacional (SIN) reduzem o isolamento do ar, causando curtos-circuitos (\textit{flashover}) e quedas de energia.

Embora o Instituto Nacional de Pesquisas Espaciais (INPE) forneça dados abertos de satélites \cite{inpe2025bdqueimadas}, cada registro bruto resume-se a uma coordenada isolada (latitude e longitude) que não informa qual comunidade ou infraestrutura está em risco. Além disso, cruzar milhões de coordenadas contra mapas vetoriais no momento da consulta torna bancos relacionais lentos. Para solucionar esse problema, desenvolveu-se na Iniciação Científica (IC) a plataforma \textbf{SIRDA 1.0} \cite{sampaio2026ic}: uma arquitetura que antecipa o cruzamento espacial em memória e agrupa focos vizinhos em perímetros de incêndio (via ST-DBSCAN) antes de transmitir os eventos via Apache Kafka para o banco colunar Apache Druid \cite{yang2014druid} e painéis no Apache Superset. Ao todo, unificaram-se 4,05 milhões de registros de 7 bases públicas (INPE, ONS, ANEEL, FUNAI, CNES/DATASUS, DNIT e IBGE), monitorando hospitais do SUS, terras indígenas, rodovias federais, municípios e a rede elétrica por bioma.

A operação da plataforma sobre um ano completo de dados validou a estabilidade da arquitetura e revelou padrões sazonais de pressão térmica com consultas em frações de segundo \cite{sampaio2026ic}. A partir dessa base consolidada, este Trabalho de Conclusão de Curso (TCC) propõe evoluir o sistema para o \textbf{SIRDA 2.0}, integrando modelos de linguagem (via protocolo FastMCP e busca documental RAG), automação de fluxos analíticos e alertas (Webhooks), processamento analítico híbrido (DuckDB \cite{raasveldt2019duckdb} e Apache Druid) e modelagem direcional de propagação do fogo para apoiar a tomada de decisão rápida baseada em dados em infraestruturas estratégicas.

\section{Motivação e Oportunidades de Evolução do Projeto}
\label{sc:motivacao}

Os bons resultados obtidos com a plataforma na Iniciação Científica motivaram a continuidade da pesquisa neste TCC, buscando enriquecer o sistema por meio de quatro frentes de aprimoramento:

\textbf{1. De Painéis de Visualização para Alertas Automáticos:} A versão atual conta com painéis interativos detalhados no Apache Superset, nos quais o usuário explora mapas e filtros na tela. Adicionar um mecanismo que envie alertas automáticos assim que um evento atingir alto risco (sem depender que alguém esteja olhando o painel naquele instante) amplia a capacidade do sistema de apoiar brigadistas e equipes da Defesa Civil em tempo hábil \cite{endsley1995toward}.

\textbf{2. Consultas em Linguagem Natural e Relatórios Prontos:} Embora os painéis visuais cubram as visões principais, investigações específicas sobre o banco de dados hoje exigem escrever consultas técnicas em SQL. A inclusão de um assistente conectado ao banco permite que um profissional faça perguntas em português e receba imediatamente um boletim padronizado em PDF, facilitando o uso da ferramenta por pessoas fora da área de computação.

\textbf{3. Inclusão do Vento e da Vegetação na Direção do Fogo:} Na versão atual, o sistema mede o risco de proximidade usando um raio igual para todos os lados (um círculo de 5 km a 15 km ao redor do fogo). Porém, ventos fortes (que chegaram a 48,2 km/h em episódios recentes segundo o ERA5) e vegetações secas de pastagem fazem o fogo correr muito mais rápido a favor do vento \cite{mapbiomas2026fogo}. Considerar a direção do vento permite esticar a área de alerta para o lado correto e reduzir alarmes desnecessários no sentido contrário.

\textbf{4. Aprofundamento Analítico com o Histórico Oficial do Setor Elétrico:} Apesar de o projeto monitorar hospitais do SUS, rodovias e terras indígenas, o setor elétrico oferece uma oportunidade especial de validação estatística por contar com o acervo público do ONS, que registra 8.395 interrupções de carga entre 2007 e 2026 (717 GW acumulados) \cite{ons2026dadosabertos}. Cruzar os dados gerados pela nossa plataforma com esse histórico oficial permite demonstrar quantitativamente como os padrões de queimadas afetam a Rede Básica ($\ge 230$ kV).

\newpage
\section{Revisão bibliográfica}
\label{sc:revisao}





A seção de revisão bibliográfica tem como objetivo apresentar e discutir os principais estudos e contribuições já realizados sobre o tema, organizando o conhecimento existente de forma lógica e crítica. Para preenchê-la:

\begin{enumerate}
    \item \textbf{Organize o conteúdo}: Estruture a revisão por temas, abordagens, cronologia ou outras categorias que facilitem a compreensão do leitor.
    \item \textbf{Sintetize os principais trabalhos}: Destaque os estudos mais relevantes, mencionando seus autores, métodos, principais achados e limitações.
    \item \textbf{Identifique lacunas}: Aponte aspectos ainda pouco explorados ou controversos, enfatizando oportunidades para contribuições futuras.
    \item \textbf{Relacione com seu tema}: Explique como os estudos revisados estão conectados ao objetivo do seu trabalho e de que forma embasam sua discussão.
    \item \textbf{Use referências de qualidade}: Priorize fontes confiáveis, como artigos científicos, livros e relatórios reconhecidos na área.
\end{enumerate}

Mantenha o texto coeso e evite apenas listar trabalhos; ao invés disso, discuta as contribuições e conexões entre eles para oferecer uma visão abrangente e crítica do estado da arte.

\section{Objetivos}
\label{sc:objetivos}

\subsection{Objetivo Geral}
Desenvolver, implementar e validar uma arquitetura de inteligência geoespacial ativa (\textit{Actionable GeoAI DataOps}) para mitigação do risco de incêndios sobre infraestruturas do SIN e áreas protegidas, integrando modelagem anisotrópica por vento, indexação Uber H3, alarmes via Webhooks e agentes de IA para relatórios situacionais.

\subsection{Objetivos Específicos}
\begin{enumerate}
    \setlength\itemsep{0em}
    \item Integrar dados horários de vento (ECMWF ERA5 / INMET) e histórico de perturbações do ONS ao pipeline Medallion;
    \item Discretizar feições e calcular o risco espacial via células hexagonais Uber H3 (\texttt{uint64}), substituindo junções geométricas por \textit{hash joins} $O(1)$;
    \item Desenvolver e calibrar o algoritmo de buffers elípticos anisotrópicos governado pelo modelo de Rothermel e classes de combustível do MapBiomas Fogo;
    \item Construir o motor de alarmes com máquina de estados de \textit{cooldown} (anti-spam) e despacho assíncrono via Webhooks (Telegram e HTTP/REST);
    \item Implementar servidor FastMCP dotado de ferramentas determinísticas (\textit{Text-to-SQL} no Apache Druid e grafo elétrico NetworkX) para geração automática de relatórios executivos no padrão ICS-209 com garantia de zero alucinação;
    \item Avaliar a acurácia espacial, a latência de triagem ($< 30\text{ s}$) e a fidelidade factual do sistema frente a métodos isotrópicos tradicionais.
\end{enumerate}

\section{Materiais e Métodos}
\label{sc:materiais_e_metodos}

\subsection{Padrão Arquitetural e Stack Tecnológico}
O sistema adota o padrão \textbf{FCIS (\textit{Functional Core, Imperative/OO Shell})}. As transformações matemáticas, sanitizações e cálculos de elipse são modelados como funções puras, sem estado e vetorizadas (\textit{DataFrame $\rightarrow$ DataFrame}). A casca orientada a objetos é restrita ao gerenciamento determinístico de conexões e soquetes (Kafka, Druid, FastMCP) via \textit{Context Managers} com injeção de dependência. O ecossistema é operado em Python 3.12 sob governança do gerenciador Astral \texttt{uv}.

A infraestrutura tecnológica integra: (i) Dagster para orquestração; (ii) Apache Kafka para mensageria; (iii) Apache Druid (OLAP colunar particionado por tempo) e DuckDB (\textit{spatial} e \textit{h3}) para persistência em GeoParquet; (iv) Apache Superset com camadas H3 (\textit{deck.gl}) e biblioteca \texttt{h3-py} para indexação em \texttt{uint64}; (v) Servidor FastMCP para interface LLM determinística, templates WeasyPrint/Markdown e despachante assíncrono de Webhooks; e (vi) Bases de dados: INPE VIIRS (375 m) e GOES-19 (10 min), reanálise horária ERA5, dados abertos ONS e MapBiomas Fogo.

\subsection{Procedimentos Metodológicos}
A metodologia estrutura-se em cinco etapas sequenciais:
\begin{enumerate}
    \setlength\itemsep{0em}
    \item \textbf{Ingestão Contratual (Bronze):} Coleta de anomalias térmicas e séries climáticas com validação via Pydantic;
    \item \textbf{Indexação H3 e Modelagem Anisotrópica (Silver):} Discretização das feições de interesse nas resoluções H3 7 ($\sim 5\text{ km}^2$) e 9 ($\sim 0,1\text{ km}^2$). Cálculo dos semieixos ($a, b$) e azimute da elipse acoplados ao vetor de vento e à carga combustível local;
    \item \textbf{Agregação e Carga Analítica (Gold):} Cruzamento das células H3 ativas contra a infraestrutura via \textit{hash joins} $O(1)$. Carga da tabela Fato no Druid no modelo analítico Snowflake (estritamente FKs numéricas, métricas e timestamps);
    \item \textbf{Inteligência Ativa (MCP e Webhooks):} Avaliação periódica de limiares críticos com máquina de \textit{cooldown} (janela de 6h). Ativação do FastMCP para executar \textit{golden queries} no Druid, disparando o boletim SitRep (ICS-209);
    \item \textbf{Validação Experimental:} Comparação de falsos positivos entre buffers circulares e elípticos, auditoria de latência de ponta a ponta ($< 30\text{ s}$) e teste de consistência factual das respostas do agente de IA.
\end{enumerate}

\section{Resultados Esperados}
\label{sc:resultados}

\subsection{Resultados Preliminares (Iniciação Científica)}
A viabilidade técnica do ecossistema encontra-se consolidada pela IC prévia: (i) pipeline Medallion funcional orquestrado por Dagster, processando mais de 4,05 milhões de registros; (ii) tabela analítica no Apache Druid respondendo consultas OLAP multidimensionais em milissegundos; (iii) grafo topológico do SIN em NetworkX para análise de criticidade $N-1$; e (iv) agrupamento espaço-temporal por ST-DBSCAN validado com sucesso.

\subsection{Resultados Esperados para o TCC02}
Com a execução do plano proposto, espera-se alcançar:
\begin{enumerate}
    \setlength\itemsep{0em}
    \item \textbf{Precisão Geométrica de Risco:} Eliminação dos alarmes falsos a barlavento característicos de buffers circulares estáticos, projetando a mancha de perigo na direção do avanço convectivo do fogo;
    \item \textbf{Aceleração Algorítmica Geoespacial:} Redução drástica da latência de cruzamento espacial em tempo real, viabilizando \textit{hash joins} em tempo constante $O(1)$ com células hexagonais Uber H3 (\texttt{uint64});
    \item \textbf{Despacho Proativo e Relatórios Instantâneos (SitRep):} Redução do tempo de triagem de incidentes de horas para menos de 30 segundos, mediante disparos automatizados de alertas via Webhooks e relatórios PDF no padrão ICS-209 com mapa vetorial anexado;
    \item \textbf{Democratização da Análise Operacional com IA:} Interface conversacional determinística baseada em FastMCP (\textit{Text-to-SQL}), permitindo que operadores consultem o banco colunar em linguagem natural sem alucinações;
    \item \textbf{Publicação Científica e Transferência:} Submissão em conferência qualificada e disponibilização do código aberto e reprodutível.
\end{enumerate}

\section{Cronograma de Atividades do TCC02}
\label{sc:cronograma}

O cronograma de execução para o TCC02 distribui-se em seis etapas sequenciais entre março e agosto de 2027:

\begin{center}
\footnotesize
\begin{ganttchart}[
    y unit title=0.42cm,
    y unit chart=0.44cm,
    x unit=1.40cm,
    hgrid,
    vgrid,
    time slot format=isodate-yearmonth,
    time slot unit=month,
    bar height=0.58,
    bar/.append style={fill=gray!45, draw=black!70, line width=0.4pt},
    bar label font=\footnotesize\bfseries,
    title label font=\footnotesize\bfseries
]{2027-03}{2027-08}
\gantttitlecalendar{year, month} \\
\ganttbar{F.1: Dados Climáticos e Elipses de Vento}{2027-03}{2027-03} \\
\ganttbar{F.2: Malha H3 e Testes em Batch (DuckDB)}{2027-03}{2027-04} \\
\ganttbar{F.3: Simulação de Streaming (Kafka / Druid)}{2027-04}{2027-05} \\
\ganttbar{F.4: Motor de Alertas e Webhooks (Cooldown)}{2027-05}{2027-06} \\
\ganttbar{F.5: Servidor FastMCP, RAG e Relatórios PDF}{2027-06}{2027-07} \\
\ganttbar{F.6: Validação Final e Escrita do TCC02}{2027-07}{2027-08}
\end{ganttchart}
\end{center}


\renewcommand\refname{Referências Bibliográficas}
\let\oldthebibliography\thebibliography
\renewcommand{\thebibliography}[1]{%
  \oldthebibliography{#1}%
  \setlength{\itemsep}{0pt}%
  \setlength{\parskip}{0pt}%
  \footnotesize
}
\bibliographystyle{plain}
\bibliography{referencias.bib}

\end{document}
